\section{Conclusion}
\label{sec:conclusion}

We have traced the behavior of a spin-$1/2$ Heisenberg pentamer chain from
isolated local multiplets to collective response.
At the tested positive couplings, the pentamer retains an even-parity spin-$1/2$
doublet, positive inter-cell exchange generates an antiferromagnetic effective
chain, and the leading projection error is quadratic in the inter-cell coupling.
Entanglement, finite-size gaps, correlations, and static structure factors
identify a $c=1$ effective chain at those points.
Reflection parity makes the lowest excited doublet dark at $q=\pi$ and shifts most
local inelastic weight into spin-$3/2$ sectors.
Correction-vector and real-time calculations connect this local selection rule to
collective spectra obtained by two independent matrix-product-state routes.

On the negative-$\alpha$ side, the isolated cell switches to a spin-$3/2$ quartet
whose response peaks at $0.60\pi$ at every sampled point, and the
accepted calculations reveal a reconstruction from $\qpeak\simeq0.60\pi$ to
$q=\pi$.
All 29 final refinement calculations converge, and a
terminal $k=8\rightarrow12$ control shows that the most sensitive $L=80$ endpoint
is stable with respect to a larger Krylov subspace.
None of the eight candidate intervals satisfies the momentum and dominant-ratio
criteria simultaneously, so the reconstruction is best described as a
finite-size crossover.

For an odd-unit-cell cluster chain, the low-energy physics is controlled not by an
inherited phase label but by which local multiplet survives, how it couples
between cells, and which of its excitations the symmetry of the probe allows one
to see.
